\addtocontents{toc}{\protect\Needspace{3\baselineskip}}
\originalchapter{玻色弦的量子谱与临界维数}\label{ch:basic-bosonic}
对应 8.251 第 18--19 讲. 量子弦需要三个相互关联的检查: 正则振子代数、正规序常数、目标洛伦兹代数. 一个负正规序常数不能只靠把发散级数“算成负数”来解释.

\originalsection{层数与正规序哈密顿量}
令 $d_\perp=D-2$, 定义
\[
 N=\sum_{n>0}\alpha_{-n}^I\alpha_n^I
 =\sum_{n>0}n a_n^{I\dagger}a_n^I,
 \qquad L_0^\perp=\alpha'p^Ip^I+N.
\]
层数是模式频率加权后的占据数, 不等于总振子个数. 因 $[N,a_n^{I\dagger}]=na_n^{I\dagger}$, 第 $n$ 模增加层数 $n$.
采用正的截距记号 $a$, 定
\[
 2\alpha'p^+p^-=L_0^\perp-a,
 \qquad M^2=\frac{N-a}{\alpha'}.
\]
原课有时以负常数 $a_{\rm notes}=-a$ 写成 $L_0+a_{\rm notes}$, 两者只是约定相反. 对世界面参数 $\tau$ 的哈密顿量是 $H_\tau=2\alpha'p^+p^-$; 对目标时间 $X^+$ 的则是 $p^-$, 不应把这两个演化生成元混同.

\originalsection{零点能的有限部分}
形式上的对称序零点能为 $E_0=(d_\perp/2)\sum_{n>0}n$. 加入指数调节 $\epsilon>0$,
\[
 \sum_{n>0}n\ee^{-\epsilon n}
 =\frac{\ee^{-\epsilon}}{(1-\ee^{-\epsilon})^2}
 =\frac1{\epsilon^2}-\frac1{12}+O(\epsilon^2).
\]
这是收敛函数的 Taylor 展开, 没有声称普通正数级数收敛到 $-1/12$. 扣除发散的局部常数后有限项为 $E_{0,\rm finite}=-d_\perp/24$, 因而该方案给 $a=d_\perp/24$. 是否允许这个扣除和是否存在额外有限移位, 还须与规范和洛伦兹一致性共同判断. 单独作 ζ 函数解析延拓得到同一有限项, 不能替代其物理检查.

\originalsection{Virasoro 代数的非中心部分}
对所有 $n$ 定义正规序 $L_n=\frac12\sum_k:\alpha_{n-k}^I\alpha_k^I:$.
利用 $[AB,C]=A[B,C]+[A,C]B$, 两项各贡献半份, 得
\[
 [L_m,\alpha_n^I]=-n\alpha_{m+n}^I.
\]
将其作用于 $L_n$ 的两个振子, 不含正规序残项的部分是 $(m-n)L_{m+n}$. 因为中心项与振子全部对易且级数守恒, 只能在 $m+n=0$ 时出现常数. 它可以在零横向动量真空上完整算出, 不需把所有无穷模式盲目相乘.

\originalsection{中心项怎样数出来}
设 $m>0$, $L_m|0\rangle=0,L_0|0\rangle=0$. 有
\[
 L_{-m}|0\rangle=\frac12\sum_{k=1}^{m-1}
 \alpha_{-(m-k)}^I\alpha_{-k}^I|0\rangle.
\]
计算范数时两个 Wick 配对各给 $k(m-k)$, 指标和给 $d_\perp$, 因此
\[
 \langle0|L_mL_{-m}|0\rangle
 =\frac{d_\perp}{2}\sum_{k=1}^{m-1}k(m-k)
 =\frac{d_\perp}{12}m(m^2-1).
\]
最后一步用 $\sum k=m(m-1)/2$, $\sum k^2=m(m-1)(2m-1)/6$. 得完整横向代数
\[
 [L_m,L_n]=(m-n)L_{m+n}
 +\frac{D-2}{12}(m^3-m)\delta_{m+n,0}.
\]
中央荷来自每个正范数横向标量的一份贡献, 不是额外人为添加的噪声. $m=0,\pm1$ 的中心项为零, 与全局共形子代数相容.

\originalsection{目标洛伦兹代数的临界条件}
光锥约束把纵向振子写成 $\alpha_n^-=L_n/(\sqrt{2\alpha'}p^+)$, 零模使用 $L_0-a$. 目标 Lorentz 荷的难点是 $M^{-I}$, 因它同时含横向振子双线性和零模非对易项. 使用厄米正规序可写
\begin{align*}
 M^{-I}={}&x_0^-p^I-\tfrac12\{x_0^I,p^-\}\\
 &-\frac{\ii}{\sqrt{2\alpha'}p^+}
 \sum_{n>0}\frac1n( L_{-n}\alpha_n^I-\alpha_{-n}^IL_n).
\end{align*}
各项的维数均为无量纲, 且花括号保证零模部分厄米. 一致的洛伦兹代数要求 $[M^{-I},M^{-J}]=0$.

公开第 19 讲给出的此长对易式的模式系数, 换成本稿正截距记号, 正比于
\[
 A_n=\frac{26-D}{24}n+
 \left(\frac{D-2}{24}-a\right)\frac1n.
\]
它来自三类项: 振子 $L$ 的普通对易产生线性 $n$ 项, 中心项产生 $n-1/n$ 组合, 零模与截距贡献 $1/n$ 项. 本稿已在前节推导中心项; 上述完整 $M^{-I}$ 双重无穷求和的逐项正规序消去仍作为原讲义长算式引用, 不将其冒称已经逐项证明. 按给出的条件对 $n=1,2$ 分别要求 $A_n=0$, 即
\[
 \frac{26-D}{24}+\frac{D-2}{24}-a=0,
 \quad2\frac{26-D}{24}+\frac12\left(\frac{D-2}{24}-a\right)=0,
\]
解得 $D=26,a=1$, 并使所有 $n$ 同时为零. 调节所得 $a=(D-2)/24$ 与此相容. 因而结论针对平直、自由、仅有玻色坐标的临界弦, 不排除有线性膨胀子或额外内部共形理论的其他构造. 完整 boost 对易式的独立算符复核列入待补高阶细节, 不能用数值验证几个低层就代替一般证明.

\originalsection{状态空间与低质量层}
在固定 $(p^+,p^I)$ 的振子真空上,
\[
 |\{k_{nI}\}\rangle=\prod_{n,I}\frac{(a_n^{I\dagger})^{k_{nI}}}{\sqrt{k_{nI}!}}
 |p^+,p^I\rangle,\quad N=\sum_{n,I}nk_{nI}.
\]
有限能层只允许有限激发. $D=26$ 时 $M^2=(N-1)/\alpha'$. $N=0$ 是标量快子; $N=1$ 有 24 个 $a_1^{I\dagger}$ 态, 对应 26 维无质量向量的横向极化. “向量态”可在附加端点标签后成为非阿贝尔规范场, 不能说自然界唯一无质量玻色子只有光子.

$N=2$ 包括 $\alpha_{-2}^I$ 的 24 态和 $\alpha_{-1}^I\alpha_{-1}^J$ 的 300 个对称态, 合计 324. 在 massive 静止系的 $SO(25)$ 小群中, 它们组成对称无迹二阶张量, 维数 $25\cdot26/2-1=324$. 这是不同光锥类型组合成完整洛伦兹表示的一次实际检验.

\originalsection{快子表示背景不稳定}
若标量有效势在 $\phi=0$ 附近为 $V=V_0-\mu^2\phi^2/2+\cdots$, 则小扰动频率满足 $\omega_k^2=k^2-\mu^2$. 对 $k<\mu$, 解为指数增长或衰减, 表明该背景不是稳定真空. 不能把它解释为可稳定发送超光速信号的普通粒子.

玻色开弦的快子与 D 膜背景的不稳定有关; 闭弦快子的背景问题更广泛. 原课介绍了开弦场论中快子凝聚与膜张力抵消的研究结果, 但仅从自由谱无法推出完整非微扰终态. 本卷不把 2007 年的研究综述当作已经在这里推导的弦场论定理.
